\newpage
\section{Topologie metrischer Räume}
\begin{definition}{Offene und abgeschlossene Mengen}
    Sei $(X,d)$ ein metrischer Raum.
    \begin{itemize}
        \item[(i)] Eine Menge $U\subset X$ heißt \emph{Umgebung} von $x\in X$, wenn ein $\varepsilon>0$ existiert mit $B_\varepsilon(x)\subset U$.
        \item[(ii)] Eine Menge $U\subset X$ heißt \emph{offen}, wenn sie Umgebung jedes ihrer Punktes $x\in U$ ist, d.h. $\forall x \in U \exists \epsilon > 0$ mit $B_\epsilon(x)\subset U$.
        \item[(iii)] Eine Menge $A\subset X$ heißt \emph{abgeschlossen}, wenn ihr Komplement $A^C := X\setminus A$ offen ist.
    \end{itemize}
\end{definition}

\begin{beispielbox}
    COMMING SOON
\end{beispielbox}

\begin{satz}{Hausdorff'sches Trennungsaxiom}
    In jedem metrischen Raum existieren zu zwei verschiedenen Punkten $x\neq y$ offene Umgebungen $U$ von $x$ und $V$ von $y$, sodass $U\cap V=\emptyset$.
\end{satz}

\begin{satz}{Eigenschaften offener Mengen}
    Sei $(X,d)$ ein metrischer Raum. Dann gilt:
    \begin{itemize}
        \item[(1)] Der Durchschnitt zweier offener Mengen ist offen.
        \item[(2)] Die Vereinigung beliebig vieler offener Mengen ist offen.
    \end{itemize}
\end{satz}

\begin{satz}{Eigenschaften abgeschlossener Mengen}
    Sei $(X,d)$ ein metrischer Raum.
    \begin{itemize}
        \item[(1)] Die Vereinigung zweier abgeschlossener Mengen ist abgeschlossen.
        \item[(2)] Der Durchschnitt beliebig vieler abgeschlossener Mengen ist abgeschlossen.
    \end{itemize}
\end{satz}

\begin{definition}{Rand einer Menge}
    Sei $(X,d)$ ein metrischer Raum und $Y\subset X$. 
    
    Ein Punkt $x\in X$ heißt \emph{Randpunkt} von $Y$, wenn jede Umgebung von $x$ sowohl Punkte aus $Y$ als auch aus $X\setminus Y$ enthält. 
    
    Die Menge aller Randpunkte heißt der \emph{Rand} $\partial Y$ von $Y$.
\end{definition}

\begin{beispielbox}
    COMMING SOON
\end{beispielbox}

\begin{definition}{Offene Kugel}
    Sei $(X,d)$ ein metrischer Raum. Für $x\in X$ und $R > 0$ ist die offene Kugel definiert durch
    $$
        B_R(x) := \{y\in X : d(x,y) < R\}.
    $$
\end{definition}

\begin{definition}{Offene und abgeschlossene Mengen}
    Sei (X,d) ein metrischer Raum.
    \begin{itemize}
        \item[(i)] $U\subset X$ heißt Umgebung von $x\in X$, wenn ein $\varepsilon>0$ existiert mit $B_\varepsilon(x)\subset U$.
        \item[(ii)] $U$ heißt offen, wenn sie Umgebung jedes ihrer Punkte ist.
        \item[(iii)] $A\subset X$ heißt abgeschlossen, wenn das Komplement $X\setminus A$ offen ist.
    \end{itemize}
\end{definition}

\begin{satz}{Hausdorff'sches Trennungsaxiom}
    In jedem metrischen Raum existieren zu zwei verschiedenen Punkten $x\neq y$ offene Umgebungen U von x und V von y, sodass $U\cap V=\emptyset$.
\end{satz}

\begin{satz}{Eigenschaften offener Mengen}
    Sei (X,d) ein metrischer Raum.
    \begin{itemize}
        \item[(1)] Der Durchschnitt zweier offener Mengen ist offen.
        \item[(2)] Die Vereinigung beliebig vieler offener Mengen ist offen.
    \end{itemize}
\end{satz}

\begin{satz}{Eigenschaften abgeschlossener Mengen}
    Sei (X,d) ein metrischer Raum.
    \begin{itemize}
        \item[(1)] Die Vereinigung zweier abgeschlossener Mengen ist abgeschlossen.
        \item[(2)] Der Durchschnitt beliebig vieler abgeschlossener Mengen ist abgeschlossen.
    \end{itemize}
\end{satz}

\begin{definition}{Rand einer Menge}
    Sei $(X,d)$ ein metrischer Raum und $Y\subset X$. Ein Punkt $x\in X$ heißt Randpunkt von $Y$, wenn jede Umgebung von $x$ sowohl Punkte aus $Y$ als auch aus $X\setminus Y$ enthält. Die Menge aller Randpunkte heißt der Rand $\partial Y$.
\end{definition}

\begin{definition}{Inneres und Abschluss}
    Sei (X,d) ein metrischer Raum und $Y\subset X$.
    \begin{itemize}
        \item[(i)] Das \emph{Innere} von $Y$ ist die Menge $Y^0 := Y \setminus \partial Y$.
        \item[(ii)] Der \emph{Abschluss} von $Y$ ist die Menge $\overline{Y}:= Y \cup \partial Y$.
    \end{itemize}
\end{definition}

\begin{satz}{Charaktersierung von Abschluss, Inneres und Rand}
    Sei $(X,d)$ ein metrischer Raum und sei $Y\subset X$.

    \begin{enumerate}%[label=(\roman*)]
        \item $Y^{\circ}$ ist die größte offene, in $Y$ enthaltene Menge.
        \item $\overline{Y}$ ist die kleinste abgeschlossene Menge, die $Y$ enthält.
        \item $\partial Y$ ist abgeschlossen.
    \end{enumerate}  
\end{satz}

\begin{satz}{Charaktersierung des Abschlusses anhand von Folgen}
    Sei \( (X,d) \) ein metrischer Raum und \(A \subset X\). Folgende Aussagen sind äquivalent:

    \begin{enumerate}
        \item[(i)] \(A\) ist abgeschlossen.
        \item[(ii)] Sei \((x_k)_{k\in\mathbb{N}}\) eine beliebige konvergente Folge in \(A\) mit \(x_k \to x\). Dann gilt: \(x \in A\).
    \end{enumerate}   
\end{satz}

\begin{satz}{Charaktersierung des Randes anhand von Folgen}
    Sei $(X,d)$ ein metrischer Raum und $A\subset X$. Dann sind die folgenden Aussagen äquivalent:

    \begin{enumerate}%[label=(\roman*)]
        \item $x\in\partial A$.
        \item Es existiert eine Folge $(x_k)_{k\in\mathbb{N}}$ in $A$ mit $x_k\to x$ und eine Folge $(x_k^{c})_{k\in\mathbb{N}}$ in $A^{c}$ mit $x_k^{c}\to x$.
    \end{enumerate}   
\end{satz}

\begin{definition}{Topologie}
    Für eine Menge $X$ heißt ein System $\mathcal{T}$ von Teilmengen $\mathcal{T} \subseteq \mathcal{P}(X)$ \emph{Topologie} auf $X$, fals folgende Eigenschaften erfüllt sind:
    \begin{enumerate}
        \item Die leere Menge und die Menge selbst sind in der Topologie enthalten:
        $$
            \emptyset, X \in \mathcal{T}
        $$

        \item Für $U,V \in \mathcal{T}$ gilt auch $U \cap V \in \mathcal{T}$

        \item Für eine beliebige Indexmenge $I$ gilt
        $$
            U_i \in \mathcal{T}, \; \forall i \in I \Rightarrow \bigcup_ {i\in I} U_i \in \mathcal{T} \, .
        $$
    \end{enumerate}
    
    Ein \emph{topologischer Raum} ist ein Paar $(X, \mathcal{T})$.

    Eine Teilmenge $U \subset X$ heißt \emph{offen}, wenn $U \in \mathcal{T}$.

    $A \subset X$ heißt \emph{abgeshclossen}, wenn ihr Komplement $A^C \in \mathcal{T}$.
\end{definition}

Eine Äquivalente Definition für den metrischen Raum ist:

\begin{definition}{Topologischer Raum}
    Ein \emph{topologischer Raum} ist ein Paar $(X, \mathcal{O})$, wobei $\mathcal{O} \subseteq \mathcal{P}(X)$ eine Familie von Teilmengen ist (offene Mengen), sodass gilt:
    
    \begin{enumerate}
        \item Die Leere Menge und die Menge $X$ sind in $\mathcal{O}$: \\
        $$\emptyset, X \in \mathcal{O}$$
        \item Endliche Schnitte offener Mengen sind offen:
        $$
            U_1, \dots, U_k \in \mathcal{O},\ k \in \mathbb{N} \quad \Rightarrow \quad \bigcap_{i=1}^k U_i \in \mathcal{O}
        $$
        \item Beliebige Vereinigungen offener Mengen sind offen:
        $$
            \{ U_i \}_{i \in I} \subseteq \mathcal{O} \quad \Rightarrow \quad \bigcup_{i \in I} U_i \in \mathcal{O}
        $$
    \end{enumerate}
    
    $\mathcal{O}$ heißt dann auch \emph{Topologie auf $X$}
\end{definition}